
\documentclass[11pt,a4paper]{article}
\usepackage[utf8]{inputenc}
\usepackage[T1]{fontenc}
\usepackage[french]{babel}
\usepackage{lmodern}
\usepackage{microtype}
\usepackage[a4paper,margin=1.7cm,headheight=15pt]{geometry}
\usepackage{amsmath,amssymb,mathtools}
\usepackage{array,tabularx,multicol}
\usepackage{enumitem}
\usepackage[most]{tcolorbox}
\usepackage{xcolor}
\usepackage{tikz}
\usepackage{pgfplots}
\pgfplotsset{compat=1.18}
\usetikzlibrary{arrows.meta,calc,patterns}
\usepackage{fancyhdr}
\usepackage{fancyvrb}
\usepackage{listings}
\usepackage{hyperref}
\usepackage{pdfpages}

\definecolor{bleufonce}{RGB}{28,55,120}
\definecolor{bleuclair}{RGB}{238,243,255}
\definecolor{grisclair}{RGB}{246,246,246}
\definecolor{tableline}{RGB}{45,45,45}
\definecolor{softgray}{RGB}{246,247,249}
\definecolor{deepblue}{RGB}{25,62,115}
\definecolor{softblue}{RGB}{235,241,250}
\definecolor{accent}{RGB}{20,82,130}

\hypersetup{colorlinks=true,linkcolor=bleufonce,urlcolor=bleufonce,pdfauthor={},pdftitle={Parcours de révision -- Fonctions -- Corrigés}}
\setlength{\parindent}{0pt}
\setlength{\parskip}{0.28em}
\renewcommand{\arraystretch}{1.12}
\setlist[itemize]{leftmargin=1.25em,itemsep=0.15em,topsep=0.2em}
\setlist[enumerate,1]{label=\textbf{\arabic*.}, leftmargin=1.05cm, itemsep=0.35em, topsep=0.25em}
\setlist[enumerate,2]{label=\textbf{\alph*.}, leftmargin=0.85cm, itemsep=0.25em, topsep=0.2em}

\pagestyle{fancy}
\fancyhf{}
\lhead{}
\rhead{}
\cfoot{\thepage}

\newcommand{\e}{\mathrm e}
\newcommand{\dd}{\,\mathrm d}
\newcommand{\vect}[1]{\overrightarrow{#1}}
\newcommand{\origine}[1]{ {\scriptsize\textcolor{bleufonce}{(site #1)}}}
\newcounter{question}
\newcounter{reponse}

\tikzset{
  tabouter/.style={draw=tableline, line width=0.85pt},
  tabline/.style={draw=tableline, line width=0.55pt},
  forbidden/.style={draw=tableline, line width=0.95pt},
  vararrow/.style={-{Latex[length=2.7mm,width=2.0mm]}, line width=0.85pt, draw=accent},
  tabtext/.style={font=\small},
  tabmath/.style={font=\small},
  labelcell/.style={font=\small\bfseries, text=deepblue},
  pointvalue/.style={font=\small\bfseries},
  endvalue/.style={font=\small},
}

\newenvironment{blocquestions}[1]{%
  \begin{tcolorbox}[
    enhanced,
    breakable,
    colback=bleuclair,
    colframe=bleufonce,
    coltitle=white,
    colbacktitle=bleufonce,
    title={\bfseries #1},
    fonttitle=\bfseries,
    boxrule=0.6pt,
    arc=2mm,
    left=2mm,
    right=2mm,
    top=2mm,
    bottom=1.5mm,
    before skip=3mm,
    after skip=3mm, fontupper=\large
  ]
  \large
  \begin{enumerate}[leftmargin=*,label=\textbf{\arabic*.},itemsep=0.75em,topsep=0.5em]
  \setcounter{enumi}{\value{question}}
}{%
  \setcounter{question}{\value{enumi}}
  \end{enumerate}
  \end{tcolorbox}
}

\newenvironment{blocreponses}[1]{%
  \begin{tcolorbox}[
    enhanced,
    breakable,
    colback=grisclair,
    colframe=bleufonce,
    coltitle=white,
    colbacktitle=bleufonce,
    title={\bfseries #1},
    fonttitle=\bfseries,
    boxrule=0.6pt,
    arc=2mm,
    left=2mm,
    right=2mm,
    top=2mm,
    bottom=1.5mm,
    before skip=3mm,
    after skip=3mm, fontupper=\large
  ]
  \large
  \begin{enumerate}[leftmargin=*,label=\textbf{\arabic*.},itemsep=0.45em,topsep=0.5em]
  \setcounter{enumi}{\value{reponse}}
}{%
  \setcounter{reponse}{\value{enumi}}
  \end{enumerate}
  \end{tcolorbox}
}

\newtcolorbox{correctionbox}[1]{
  enhanced,
  breakable,
  colback=grisclair,
  colframe=black!55,
  boxrule=0.55pt,
  arc=2mm,
  left=2mm,right=2mm,top=1.5mm,bottom=1.5mm,
  title=\textbf{#1},
  coltitle=black,
  colbacktitle=black!12,
  before skip=2mm,
  after skip=2mm
}

\newcommand{\titreSujet}[2]{%
  \subsection*{#1}
  \addcontentsline{toc}{subsection}{#1}
  \textit{Référence : #2}\par\medskip\hrule\medskip
}

\lstset{language=Python,basicstyle=\ttfamily\small,numbers=left,numberstyle=\tiny,frame=single,showstringspaces=false,columns=fullflexible,keepspaces=true}
\sloppy

\begin{document}
\begin{center}
{\LARGE\bfseries Parcours de révision -- Fonctions}\\[1mm]

\end{center}
\vspace{1mm}
\tableofcontents
\clearpage

\section*{Partie 1 -- Corrections des questions flashs}
\addcontentsline{toc}{section}{Partie 1 -- Corrections des questions flashs}
\setcounter{reponse}{0}
\begin{blocreponses}{Réponses -- A. Vocabulaire essentiel}
  \item \origine{266} Une fonction de la forme $x\mapsto ax+b$.
  \item \origine{267} Une fonction de la forme $x\mapsto ax^2+bx+c$, avec $a\neq0$.
  \item \origine{268} Une fonction écrite comme quotient de deux polynômes.
  \item \origine{269} L'ensemble des réels pour lesquels la fonction a un sens.
  \item \origine{270} Une valeur exclue du domaine de définition.
  \item \origine{274} Un réel $\alpha$ tel que $f(\alpha)=0$.
  \item \origine{275} Une fonction dérivable en tout point de cet intervalle.
  \item \origine{276} Le nombre $f'(a)$ au point d'abscisse $a$.
  \item \origine{277} Un maximum local ou un minimum local.
  \item \origine{280} Une fonction qui y est toujours croissante ou toujours décroissante.
\end{blocreponses}
\begin{blocreponses}{Réponses -- B. Limites, continuité et TVI}
  \item \origine{96} $+\infty$.
  \item \origine{97} $2$.
  \item \origine{98} $4$.
  \item \origine{99} Une forme pour laquelle on ne peut pas conclure directement sur la limite.
  \item \origine{103} Le terme de plus haut degré.
  \item \origine{104} Une droite $x=a$ lorsque $f(x)\to\pm\infty$ quand $x\to a$.
  \item \origine{105} Une droite $y=\ell$ lorsque $f(x)\to\ell$ quand $x\to+\infty$.
  \item \origine{107} $\displaystyle\lim_{x\to a} f(x)=f(a)$.
  \item \origine{109} Le théorème des valeurs intermédiaires.
  \item \origine{110} Que la fonction soit strictement monotone.
\end{blocreponses}
\begin{blocreponses}{Réponses -- C. Exponentielle et logarithme}
  \item \origine{111} $\ln$.
  \item \origine{112} Le logarithme népérien.
  \item \origine{113} Sur $]0,+\infty[$.
  \item \origine{114} $0$.
  \item \origine{115} $1$.
  \item \origine{116} $1$.
  \item \origine{117} $x$.
  \item \origine{118} $x$.
  \item \origine{119} $\ln(ab)=\ln(a)+\ln(b)$.
  \item \origine{120} $\ln\left(\dfrac{a}{b}\right)=\ln(a)-\ln(b)$.
\end{blocreponses}
\begin{blocreponses}{Réponses -- D. Dérivation, tangente, variations et convexité}
  \item \origine{125} $e^x$.
  \item \origine{126} $a e^{ax+b}$.
  \item \origine{127} $\dfrac1x$.
  \item \origine{128} $\dfrac{a}{ax+b}$, sur l'intervalle où cela a un sens.
  \item \origine{139} Le coefficient directeur de la tangente à la courbe au point d'abscisse $a$.
  \item \origine{140} $y=f'(a)(x-a)+f(a)$.
  \item \origine{141} La tangente est horizontale au point d'abscisse $a$.
  \item \origine{142} Les variations de $f$.
  \item \origine{146} Une fonction dont la courbe est au-dessus de ses tangentes.
  \item \origine{147} Le signe de $f''$.
\end{blocreponses}
\begin{blocreponses}{Réponses -- E. Primitives, intégrales, équations différentielles et calculs flash}
  \item \origine{155} Une fonction $F$ telle que $F'=f$ sur cet intervalle.
  \item \origine{157} Une primitive est $e^x$.
  \item \origine{160} Une primitive est $\ln(x)$.
  \item \origine{164} Si $F'=f$, alors $\displaystyle\int_a^b f(x)\,dx=F(b)-F(a)$.
  \item \origine{166} Une équation différentielle sans second membre.
  \item \origine{167} $y(x)=Ce^{-ax}$, avec $C\in\mathbb R$.
  \item \origine{237} $x\mapsto 6x$.
\end{blocreponses}
\clearpage
\section*{Partie 2 -- Corrigés des sujets de bac}
\addcontentsline{toc}{section}{Partie 2 -- Corrigés des sujets de bac}
\titreSujet{Sujet 1 -- Exponentielle, équation différentielle, convexité et aire}{Asie, 5 septembre 2025, sujet 1, exercice 1}

\begin{correctionbox}{Affirmation 1}
Pour tout réel $x$, on a
\[
f'(x)=1\times \e^{-2x}+x\times(-2)\e^{-2x}=(1-2x)\e^{-2x}.
\]
L'affirmation est \textbf{vraie}.
\end{correctionbox}

\begin{correctionbox}{Affirmation 2}
On calcule
\[
f'(x)+2f(x)=(1-2x)\e^{-2x}+2x\e^{-2x}=\e^{-2x}.
\]
La fonction $f$ est donc solution de $y'+2y=\e^{-2x}$. L'affirmation est \textbf{vraie}.
\end{correctionbox}

\begin{correctionbox}{Affirmation 3}
On dérive $f'(x)=(1-2x)\e^{-2x}$ :
\[
f''(x)=-2\e^{-2x}-2(1-2x)\e^{-2x}=4(x-1)\e^{-2x}.
\]
Comme $\e^{-2x}>0$, le signe de $f''(x)$ est celui de $x-1$. La fonction $f$ est donc concave sur $]-\infty;1]$ et convexe sur $[1;+\infty[$. L'affirmation est \textbf{fausse}.
\end{correctionbox}

\begin{correctionbox}{Affirmation 4}
On sait que $f'(x)=(1-2x)\e^{-2x}$. Comme $\e^{-2x}>0$, $f$ est croissante sur $]-\infty;\frac12]$ puis décroissante sur $[\frac12;+\infty[$.
De plus
\[
\lim_{x\to-\infty}x\e^{-2x}=-\infty,
\qquad f\left(\dfrac12\right)=\dfrac{1}{2e}>0,
\qquad \lim_{x\to+\infty}x\e^{-2x}=0^+.
\]
Sur $]-\infty;\frac12]$, la fonction passe une seule fois de $-\infty$ à une valeur positive. Sur $[\frac12;+\infty[$, elle reste positive. L'équation $f(x)=-1$ admet donc une unique solution. L'affirmation est \textbf{vraie}.
\end{correctionbox}

\begin{correctionbox}{Affirmation 5}
Sur $[0;1]$, $f(x)=x\e^{-2x}\geqslant0$. L'aire demandée vaut donc
\[
\int_0^1 x\e^{-2x}\,dx.
\]
Une primitive est
\[
F(x)=-\left(\dfrac{x}{2}+\dfrac14\right)\e^{-2x}.
\]
Donc
\[
\int_0^1 x\e^{-2x}\,dx
=F(1)-F(0)
=-\dfrac34\e^{-2}+\dfrac14
=\dfrac14-\dfrac{3\e^{-2}}4.
\]
L'affirmation est \textbf{vraie}.
\end{correctionbox}

\newpage
\titreSujet{Sujet 2 -- Logarithme, variations et convexité}{Polynésie, 18 juin 2025, sujet 2, exercice 2}

\begin{correctionbox}{1. Conjectures graphiques}
À la lecture du graphique, on peut conjecturer que la fonction $f$ est strictement croissante sur $]2;+\infty[$, que
\[
\lim_{x\to2^+}f(x)=-\infty
\quad\text{et}\quad
\lim_{x\to+\infty}f(x)=+\infty.
\]
La droite d'équation $x=2$ semble être une asymptote verticale.
\end{correctionbox}

\begin{correctionbox}{2. Équation $f(x)=0$}
Sur $]2;+\infty[$, on a $x\neq0$, donc
\[
f(x)=0 \Longleftrightarrow x\ln(x-2)=0
\Longleftrightarrow \ln(x-2)=0
\Longleftrightarrow x-2=1.
\]
Ainsi
\[
\boxed{x=3}.
\]
\end{correctionbox}

\begin{correctionbox}{3. Limite en $2$}
Lorsque $x\to2^+$, on a $x\to2$ et $\ln(x-2)\to-\infty$. Donc
\[
\lim_{x\to2^+}x\ln(x-2)=-\infty.
\]
Cela confirme l'existence d'une asymptote verticale d'équation $x=2$.
\end{correctionbox}

\begin{correctionbox}{4. Dérivée}
Pour tout $x>2$,
\[
f'(x)=1\times\ln(x-2)+x\times\dfrac1{x-2}
=\ln(x-2)+\dfrac{x}{x-2}.
\]
\end{correctionbox}

\begin{correctionbox}{5. Étude de $g=f'$}
\begin{enumerate}
\item Pour tout $x>2$,
\[
g'(x)=\dfrac1{x-2}+\dfrac{(x-2)-x}{(x-2)^2}
=\dfrac{x-2-2}{(x-2)^2}
=\dfrac{x-4}{(x-2)^2}.
\]

\item Comme $(x-2)^2>0$, le signe de $g'(x)$ est celui de $x-4$. Donc $g$ est décroissante sur $]2;4]$ puis croissante sur $[4;+\infty[$.
De plus
\[
g(4)=\ln2+\dfrac42=\ln2+2.
\]

\begin{center}
\begin{tikzpicture}[x=1cm,y=1cm]
\def\L{2.1}\def\W{10.4}\def\H{2.55}\def\Yx{1.75}\def\Xm{5.8}
\fill[softgray] (0,\Yx) rectangle (\W,\H);
\fill[softblue] (0,0) rectangle (\L,\Yx);
\draw[tabouter] (0,0) rectangle (\W,\H);
\draw[tabline] (0,\Yx) -- (\W,\Yx);
\draw[tabline] (\L,0) -- (\L,\H);
\draw[tabline] (\Xm,0) -- (\Xm,\Yx);
\draw[forbidden] (\L+0.06,0) -- (\L+0.06,\H);
\draw[forbidden] (\L+0.14,0) -- (\L+0.14,\H);
\node[labelcell] at ({\L/2},{(\Yx+\H)/2}) {$x$};
\node[labelcell] at ({\L/2},{\Yx/2}) {$g(x)$};
\node[pointvalue] at ({\L+0.35},{(\Yx+\H)/2}) {$2$};
\node[pointvalue] at (\Xm,{(\Yx+\H)/2}) {$4$};
\node[tabmath] at ({\W-0.6},{(\Yx+\H)/2}) {$+\infty$};
\node[endvalue] at ({\L+0.9},1.35) {$+\infty$};
\node[pointvalue] at (\Xm,0.35) {$\ln2+2$};
\node[endvalue] at ({\W-0.8},1.35) {$+\infty$};
\draw[vararrow] ({\L+1.4},1.25) -- ({\Xm-0.65},0.5);
\draw[vararrow] ({\Xm+0.65},0.5) -- ({\W-1.25},1.25);
\end{tikzpicture}
\end{center}

\item Le minimum de $g$ est $\ln2+2>0$. Donc, pour tout $x>2$, $g(x)>0$.

\item Comme $f'(x)=g(x)>0$, la fonction $f$ est strictement croissante sur $]2;+\infty[$.
\end{enumerate}
\end{correctionbox}

\begin{correctionbox}{6. Convexité}
Comme $f''(x)=g'(x)=\dfrac{x-4}{(x-2)^2}$, la fonction $f$ est concave sur $]2;4]$ et convexe sur $[4;+\infty[$.
La courbe admet un point d'inflexion pour $x=4$, de coordonnées
\[
\boxed{(4;4\ln2)}.
\]
\end{correctionbox}

\begin{correctionbox}{7. Tangentes de coefficient directeur 3}
Une tangente de coefficient directeur $3$ correspond à une solution de
\[
f'(x)=3,
\]
c'est-à-dire $g(x)=3$. Or $g$ tend vers $+\infty$ en $2^+$, décroît jusqu'à $\ln2+2<3$, puis croît vers $+\infty$.
Il y a donc deux solutions. La courbe admet donc \textbf{deux} tangentes de coefficient directeur $3$.
\end{correctionbox}

\newpage
\titreSujet{Sujet 3 -- Équation différentielle et étude d'une exponentielle}{Amérique du Sud, 21 novembre 2024, jour 1, exercice 1}

\begin{correctionbox}{Partie A}
\begin{enumerate}
\item Pour $g(x)=ax\e^{-x/4}$,
\[
g'(x)=a\e^{-x/4}-\dfrac a4x\e^{-x/4}.
\]
Donc
\[
g'(x)+\dfrac14g(x)=a\e^{-x/4}.
\]
Pour que $g$ soit solution de $(E)$, il faut $a=20$.

\item Les solutions de $y'+\frac14y=0$ sont
\[
y(x)=C\e^{-x/4},\qquad C\in\mathbb R.
\]

\item Les solutions de $(E)$ sont donc
\[
y(x)=20x\e^{-x/4}+C\e^{-x/4}=(20x+C)\e^{-x/4},\qquad C\in\mathbb R.
\]

\item La condition $f(0)=8$ donne $C=8$. Ainsi
\[
f(x)=(20x+8)\e^{-x/4}.
\]
\end{enumerate}
\end{correctionbox}

\begin{correctionbox}{Partie B}
\begin{enumerate}
\item On dérive :
\[
f'(x)=20\e^{-x/4}+(20x+8)\left(-\dfrac14\right)\e^{-x/4}
=(18-5x)\e^{-x/4}.
\]
Comme $\e^{-x/4}>0$, le signe de $f'(x)$ est celui de $18-5x$.

\begin{center}
\begin{tikzpicture}[x=1cm,y=1cm]
\def\L{2.1}\def\W{10.8}\def\H{2.75}\def\Yx{1.92}\def\Xm{5.7}
\fill[softgray] (0,\Yx) rectangle (\W,\H);
\fill[softblue] (0,0) rectangle (\L,\Yx);
\draw[tabouter] (0,0) rectangle (\W,\H);
\draw[tabline] (0,\Yx) -- (\W,\Yx);
\draw[tabline] (\L,0) -- (\L,\H);
\draw[tabline] (\Xm,0) -- (\Xm,\Yx);
\node[labelcell] at ({\L/2},{(\Yx+\H)/2}) {$x$};
\node[labelcell] at ({\L/2},{\Yx/2}) {$f(x)$};
\node[pointvalue] at ({\L+0.35},{(\Yx+\H)/2}) {$0$};
\node[pointvalue] at (\Xm,{(\Yx+\H)/2}) {$\frac{18}{5}$};
\node[tabmath] at ({\W-0.6},{(\Yx+\H)/2}) {$+\infty$};
\node[endvalue] at ({\L+0.5},0.35) {$8$};
\node[pointvalue] at (\Xm,1.45) {$80\e^{-9/10}$};
\node[endvalue] at ({\W-0.7},0.35) {$0$};
\draw[vararrow] ({\L+0.95},0.5) -- ({\Xm-0.75},1.35);
\draw[vararrow] ({\Xm+0.75},1.35) -- ({\W-1.15},0.5);
\end{tikzpicture}
\end{center}

Le maximum est atteint en $x=\frac{18}{5}$ et vaut
\[
f\left(\dfrac{18}{5}\right)=80\e^{-9/10}.
\]

\item Sur $[14;15]$, la fonction $f$ est strictement décroissante car $14>\frac{18}{5}$. De plus
\[
f(14)\approx8{,}70>8,
\qquad
f(15)\approx7{,}24<8.
\]
Par continuité et stricte monotonie, l'équation $f(x)=8$ admet une unique solution $\alpha$ dans $[14;15]$.

Le tableau d'exécution est :
\[
\begin{array}{|c|c|c|c|c|}
\hline
 a&14&14&14{,}25&14{,}375\\ \hline
 b&15&14{,}5&14{,}5&14{,}5\\ \hline
 b-a&1&0{,}5&0{,}25&0{,}125\\ \hline
 m&14{,}5&14{,}25&14{,}375&14{,}4375\\ \hline
 f(m)>8&\text{FAUX}&\text{VRAI}&\text{VRAI}&\text{VRAI}\\ \hline
\end{array}
\]
À la fin, l'algorithme renvoie l'intervalle
\[
[14{,}4375;14{,}5].
\]
Son objectif est de donner un encadrement de la solution $\alpha$ de l'équation $f(x)=8$ avec une amplitude inférieure ou égale à $0{,}1$.
\end{enumerate}
\end{correctionbox}

\newpage
\titreSujet{Sujet 4 -- Logarithme, continuité et suite associée}{Centres étrangers, 12 juin 2025, sujet 1, exercice 3}

\begin{correctionbox}{Partie A}
\begin{enumerate}
\item Lorsque $x\to-1^+$, on a $x+1\to0^+$, donc $\ln(x+1)\to-\infty$. Ainsi
\[
\lim_{x\to-1^+}f(x)=-\infty.
\]

\item Pour tout $x>-1$,
\[
f'(x)=\dfrac4{x+1}-\dfrac{2x}{25}
=\dfrac{100-2x(x+1)}{25(x+1)}
=\dfrac{100-2x-2x^2}{25(x+1)}.
\]

\item Sur $]-1;+\infty[$, le dénominateur est positif. Le signe de $f'(x)$ est donc celui de
\[
100-2x-2x^2=-2(x^2+x-50).
\]
L'unique racine appartenant à l'intervalle est
\[
\beta=\dfrac{-1+\sqrt{201}}2\approx6{,}59.
\]
La fonction $f$ est donc croissante sur $]-1;\beta]$ puis décroissante sur $[\beta;+\infty[$. Comme $6{,}5<\beta$, elle est strictement croissante sur $[2;6{,}5]$.

\item La fonction $h$ est continue sur $[2;6{,}5]$. On a
\[
h(2)=4\ln3-\dfrac4{25}-2\approx2{,}23>0,
\]
et
\[
h(6{,}5)=4\ln(7{,}5)-\dfrac{6{,}5^2}{25}-6{,}5\approx-0{,}13<0.
\]
D'après le tableau de variations fourni, $h$ est d'abord croissante puis décroissante, avec un maximum positif. Elle ne peut donc pas s'annuler sur la partie croissante, puis s'annule une seule fois sur la partie décroissante. L'équation $h(x)=0$ admet donc une unique solution $\alpha\in[2;6{,}5]$.

\item L'appel \texttt{bornes(2)} renvoie
\[
\boxed{(6{,}36;6{,}37)}.
\]
Ces valeurs donnent un encadrement au centième de la solution $\alpha$ de l'équation $f(x)=x$ :
\[
6{,}36<\alpha<6{,}37.
\]
\end{enumerate}
\end{correctionbox}

\begin{correctionbox}{Partie B}
\begin{enumerate}
\item Montrons par récurrence que $2\leqslant u_n\leqslant u_{n+1}<6{,}5$.
Pour $n=0$, $u_0=2$ et
\[
u_1=f(2)=4\ln3-\dfrac4{25}\approx4{,}23,
\]
donc $2\leqslant u_0\leqslant u_1<6{,}5$.
Supposons $2\leqslant u_n\leqslant u_{n+1}<6{,}5$. Comme $f$ est croissante sur $[2;6{,}5]$,
\[
f(2)\leqslant f(u_n)\leqslant f(u_{n+1})<f(6{,}5).
\]
Or $f(u_n)=u_{n+1}$, $f(u_{n+1})=u_{n+2}$ et $f(6{,}5)\approx6{,}37<6{,}5$. Donc
\[
2\leqslant u_{n+1}\leqslant u_{n+2}<6{,}5.
\]
La propriété est démontrée.

\item La suite $(u_n)$ est croissante et majorée par $6{,}5$, donc elle converge vers une limite $\ell$.

\item Comme $f$ est continue sur $[2;6{,}5]$, on peut passer à la limite dans $u_{n+1}=f(u_n)$ :
\[
\ell=f(\ell).
\]
Donc $h(\ell)=f(\ell)-\ell=0$. Or $\alpha$ est l'unique solution de $h(x)=0$ sur $[2;6{,}5]$, donc
\[
\ell=\alpha.
\]
\end{enumerate}
\end{correctionbox}

\newpage
\titreSujet{Sujet 5 -- Trigonométrie, convexité et intégration}{Amérique du Nord, 22 mai 2025, sujet 2, exercice 4}

\begin{correctionbox}{Partie A}
\begin{enumerate}
\item On dérive :
\[
f'(x)=\e^x\sin x+\e^x\cos x=\e^x(\sin x+\cos x).
\]
Sur $\left[0;\frac\pi2\right]$, $\sin x\geqslant0$ et $\cos x\geqslant0$, et leur somme est strictement positive. Donc $f'(x)>0$ et $f$ est strictement croissante sur cet intervalle.

\item On a $f(0)=0$ et $f'(0)=1$. La tangente en $0$ a donc pour équation
\[
y=x.
\]
Ensuite
\[
f''(x)=\e^x(\sin x+\cos x)+\e^x(\cos x-\sin x)=2\e^x\cos x.
\]
Sur $\left[0;\frac\pi2\right]$, $f''(x)\geqslant0$, donc $f$ est convexe.
Une fonction convexe est au-dessus de ses tangentes ; en particulier, au-dessus de la tangente en $0$. Donc
\[
\e^x\sin x\geqslant x
\]
pour tout $x\in\left[0;\frac\pi2\right]$.

\item Le signe de $f''(x)=2\e^x\cos x$ change lorsque $x$ traverse $\frac\pi2$ : positif avant, négatif après. Le point d'abscisse $\frac\pi2$ est donc un point d'inflexion.
\end{enumerate}
\end{correctionbox}

\begin{correctionbox}{Partie B}
\begin{enumerate}
\item Première intégration par parties, avec $u=\sin x$ et $v'=\e^x$ :
\[
I=\left[\e^x\sin x\right]_0^{\pi/2}-\int_0^{\pi/2}\e^x\cos x\,dx
=\e^{\pi/2}-J.
\]
Deuxième intégration par parties, avec $u=\e^x$ et $v'=\sin x$ :
\[
I=\left[-\e^x\cos x\right]_0^{\pi/2}+\int_0^{\pi/2}\e^x\cos x\,dx
=1+J.
\]

\item En additionnant les deux expressions,
\[
2I=1+\e^{\pi/2},
\]
donc
\[
I=\dfrac{1+\e^{\pi/2}}2.
\]

\item Sur $\left[0;\frac\pi2\right]$, on a $f(x)\geqslant x$. L'aire vaut donc
\[
\int_0^{\pi/2}\left(\e^x\sin x-x\right)\,dx
=I-\left[\dfrac{x^2}{2}\right]_0^{\pi/2}.
\]
Ainsi
\[
\mathcal A=\dfrac{1+\e^{\pi/2}}2-\dfrac{\pi^2}{8}.
\]
\end{enumerate}
\end{correctionbox}

\newpage
\titreSujet{Sujet 6 -- Équation différentielle, vitesse et intégration}{Métropole, 18 juin 2025, sujet 2, exercice 4}

\begin{correctionbox}{Partie A}
D'après le graphique :
\begin{enumerate}
\item Le chariot a parcouru environ $15$ m au bout de $2$ secondes.
\item L'asymptote horizontale est proche de $23$ m. Une longueur minimale de $23$ m semble donc suffisante.
\item Le coefficient directeur de la tangente en $A$ est environ
\[
d'(4{,}7)\approx0{,}96.
\]
Cela signifie qu'au bout de $4{,}7$ secondes, la vitesse du chariot est d'environ $0{,}96$ m.s$^{-1}$.
\end{enumerate}
\end{correctionbox}

\begin{correctionbox}{Partie B}
\begin{enumerate}
\item Les solutions de $(E')$ sont
\[
y(t)=C\e^{-0{,}6t},\qquad C\in\mathbb R.
\]
Pour $g(t)=t\e^{-0{,}6t}$,
\[
g'(t)=\e^{-0{,}6t}-0{,}6t\e^{-0{,}6t},
\]
donc
\[
g'(t)+0{,}6g(t)=\e^{-0{,}6t}.
\]
Ainsi $g$ est une solution particulière de $(E)$.
Les solutions de $(E)$ sont donc
\[
y(t)=(t+C)\e^{-0{,}6t}.
\]
Comme $v(0)=12$, on obtient $C=12$, donc
\[
v(t)=(12+t)\e^{-0{,}6t}.
\]

\item On dérive :
\[
v'(t)=\e^{-0{,}6t}-0{,}6(12+t)\e^{-0{,}6t}=(-6{,}2-0{,}6t)\e^{-0{,}6t}.
\]
Puisque $\e^{-0{,}6t}>0$ et $-6{,}2-0{,}6t<0$ pour tout $t\geqslant0$, la fonction $v$ est strictement décroissante sur $[0;+\infty[$.
De plus
\[
\lim_{t\to+\infty}v(t)=0.
\]

\begin{center}
\begin{tikzpicture}[x=1cm,y=1cm]
\def\L{2.1}\def\W{9.4}\def\H{2.55}\def\Yx{1.75}
\fill[softgray] (0,\Yx) rectangle (\W,\H);
\fill[softblue] (0,0) rectangle (\L,\Yx);
\draw[tabouter] (0,0) rectangle (\W,\H);
\draw[tabline] (0,\Yx) -- (\W,\Yx);
\draw[tabline] (\L,0) -- (\L,\H);
\node[labelcell] at ({\L/2},{(\Yx+\H)/2}) {$t$};
\node[labelcell] at ({\L/2},{\Yx/2}) {$v(t)$};
\node[pointvalue] at ({\L+0.35},{(\Yx+\H)/2}) {$0$};
\node[tabmath] at ({\W-0.6},{(\Yx+\H)/2}) {$+\infty$};
\node[pointvalue] at ({\L+0.5},1.35) {$12$};
\node[endvalue] at ({\W-0.65},0.35) {$0$};
\draw[vararrow] ({\L+0.95},1.25) -- ({\W-1.1},0.5);
\end{tikzpicture}
\end{center}

Par continuité et stricte décroissance, comme $v(0)=12>1$ et $\lim v=0<1$, l'équation $v(t)=1$ admet une unique solution $\alpha$. À la calculatrice,
\[
\alpha\approx4{,}7
\]
au dixième.

\item Le système mécanique se déclenche lorsque $v(t)\leqslant1$. La vitesse étant décroissante, il se déclenche à partir de $t=\alpha$, soit au bout d'environ
\[
\boxed{4{,}7\text{ s}}.
\]
\end{enumerate}
\end{correctionbox}

\begin{correctionbox}{Partie C}
\begin{enumerate}
\item On calcule
\[
d(t)=\int_0^t (12+x)\e^{-0{,}6x}\,dx.
\]
Une intégration par parties sur le terme $\int x\e^{-0{,}6x}\,dx$ donne
\[
d(t)=\e^{-0{,}6t}\left(-\dfrac53t-\dfrac{205}{9}\right)+\dfrac{205}{9}.
\]

\item Le dispositif se déclenche pour $t=\alpha\approx4{,}69$. Ainsi
\[
d(\alpha)=\e^{-0{,}6\alpha}\left(-\dfrac53\alpha-\dfrac{205}{9}\right)+\dfrac{205}{9}
\approx20{,}94.
\]
Selon ce modèle, le chariot parcourt environ
\[
\boxed{20{,}94\text{ m}}
\]
avant le déclenchement du dispositif.
\end{enumerate}
\end{correctionbox}
\end{document}
